\section{Past Actions and Current Situation}

\lead{The Baltic Sea is now the most closely watched stretch of water in the NATO area. With Finland and Sweden's accession, nine of the eleven coastal states are Alliance members, leading many observers to call it a “NATO lake”.\footnote{Braw, E. (2025). How the Baltic Sea nations have tackled suspicious cable cuts. Atlantic Council. \url{https://www.atlanticcouncil.org/in-depth-research-reports/issue-brief/how-the-baltic-sea-nations-have-tackled-suspicious-cable-cuts/}} However, that near encirclement has not made the sea safer.}

Beneath it run more than thirty cross-border cables and pipelines. Since 2022 about ten of them have been cut or damaged, with seven of those cuts concentrated in the three months between November 2024 and January 2025.\footnote{Buchholz, K. (2025). 11 Baltic cables damaged in 15 months, pushing NATO to boost security. DGfiGnsG KGws. \url{https://www.defensenews.com/global/europe/2025/01/28/11-baltic-cables-damaged-in-15-months-pushing-nato-to-boost-security/};} The basin has become, in the framing of one recent analysis, “ground zero for the targeting of critical undersea infrastructure”.\footnote{Jackson School of International Studies. (2025). Baltic Sea undersea cable security. University of Washington. \url{https://jsis.washington.edu/news/baltic-sea-undersea-cable-security/}}

The pattern began with Nord Stream. In September 2022 the two pipelines were ruptured by explosions in the Swedish and Danish economic zones.\footnote{Security Council Report. (2025, August). The Nord Stream incident: Open briefing. \url{https://www.securitycouncilreport.org/whatsinblue/2025/08/the-nord-stream-incident-op} en-briefing-2.php} Although Sweden recovered traces of explosives confirming deliberate sabotage, Sweden and Denmark closed their investigations in 2024 without naming a state.\footnote{Sweden ends Nord Stream sabotage probe, hands evidence to Germany. (2024, February 7). RGuťGrs. \url{https://www.reuters.com/world/europe/sweden-ends-investigation-into-nord-stream-pipe} line-blasts-2024-02-07/}

A year later, the method changed. In October 2023, the Balticconnector gas pipeline between Finland and Estonia was damaged when the Hong Kong-flagged, Chinese-owned container ship Newnew Polar Bear dragged its anchor along the seabed, striking the pipeline and nearby data cables.\footnote{Beyond the Eagle S: How Yi Peng 3 and Newnew Polar Bear wreaked havoc in Baltic Sea. \url{https://www.cbsnews.com/news/eagle-s-how-yi-peng-3-and-newnew-polar-bear-wreake} d-havoc-in-baltic-sea-60-minutes/} The most consequential case came weeks later, when the shadow fleet tanker Eagle S damaged the cables connecting Estonia and Finland on 25 December 2024, cutting the Estlink 2 power link, and Finland responded by seizing the vessel and bringing charges against its crew.\footnote{Finland boards oil tanker suspected of causing internet, power cable damage in Baltic. (2024, December 26). RGuťGrs. \url{https://www.reuters.com/world/europe/finland-police-investigate-role-foreign-ship-after-power-cable-outage-2024-12-26/}}

The incidents did not stop in 2025.

\begin{itemize}
  \item \textbf{January:} A Latvia to Sweden cable near Gotland was damaged
  \item \textbf{February:} The C-Lion1 line was hit a third time.
  \item \textbf{May:} The Estonian Navy intercepted the suspected Russian-affiliated tanker \emph{Jaguar}, operating under a Gabon flag, but could not forcibly board it given its limited naval capacity and the presence of a Russian fighter jet escorting the ship.
  \item \textbf{December:} A fiber-optic line connecting Helsinki and Tallinn was disrupted, triggering a full national response from Finnish authorities.
\end{itemize}

\subsection{NATO’s response: Baltic Sentry}

NATO's institutional response has taken shape quickly. At a Summit of Baltic Sea Allies in Helsinki on January 14, 2025, Secretary General Mark Rutte announced the launch of Baltic Sentry to enhance NATO's military presence in the Baltic Sea and improve Allies' ability to respond to destabilizing acts.\footnote{NATO launches "Baltic Sentry" to increase critical infrastructure security. (2025, January 14). North Atlantic Treaty Organization. \url{https://www.nato.int/en/news-and-events/articles/news/2025/01/14/nato-launches-baltic-sentry-to-increase-critical-infrastructure-security}}

This initiative involves:

\begin{itemize}
  \item Frigates and maritime patrol aircraft
  \item A small fleet of naval drones
  \item Integration of national surveillance systems
  \item A new NATO Maritime Centre for the Security of Critical Undersea Infrastructure
  \item A Critical Undersea Infrastructure Network that brings industry into the response
\end{itemize}

In parallel, the United Kingdom launched a Joint Expeditionary Force operation, Nordic Warden, which uses artificial intelligence to register shadow fleet vessels and signal allies when a risk is detected.\footnote{British-led force to use AI in tracking Russia's shadow fleet. (2025, January 7). DGfiGnsG KGws. \url{https://www.defensenews.com/global/europe/2025/01/07/british-led-force-to-use-ai-in-tr} acking-russias-shadow-fleet/} The system monitors 22 areas of interest, including parts of the English Channel, the North Sea, Kattegat, and Baltic Sea, from the operational headquarters in Northwood, London.\footnote{UK-led JEF activates Nordic Warden. (2025, January 6). Joint Forces News. \url{https://www.joint-forces.com/world-news/defence-news/78866-uk-led-jef-activates-nord} ic-warden} Rutte was explicit that enforcement would be firm, citing Finland's seizures as the benchmark and warning that ship captains face possible boarding, impounding, and arrest.\footnote{NATO launches "Baltic Sentry," 2025; Nato flotilla assembles off Estonia to protect undersea cables in Baltic Sea. (2025, January 19). ThG Guardian. \url{https://www.theguardian.com/world/2025/jan/19/nato-flotilla-assembles-off-estonia-prote} ct-undersea-cables-baltic-sea}

\subsection{Legal challenges}

The harder problem, and the one this committee must sit with, is that these acts are engineered to fall below the line at which NATO can clearly act:

\begin{itemize}
  \item \textbf{Freedom of navigation:} vessels in international waters enjoy freedom of navigation under UNCLOS
  \item \textbf{Port-state limitations:} port states can deny entry or inspect, but shadow fleet vessels typically avoid NATO ports,
  \item \textbf{Sanctions:} sanctions target ownership and insurance rather than the vessels themselves.
  \item \textbf{Weak treaty provisions:} UNCLOS Article 113 only obliges states to criminalise damage to a submarine cable and grants no authority to enforce against foreign ships.
  \item \textbf{Attribution:} an anchor cut is as plausibly an accident as sabotage, and that ambiguity is the point.
\end{itemize}

NATO's 2022 Strategic Concept acknowledges that hybrid operations could reach the level of an armed attack and trigger collective defence.\footnote{North Atlantic Treaty Organization. (2022). KATO 2022 SťraťGgic ConcGpť. \url{https://www.nato.int/content/dam/nato/webready/documents/publications-and-reports/str} ategic-concepts/2022/290622-strategic-concept.pdf} Yet no precedent exists for infrastructure sabotage crossing that threshold, and the only Article 5 invocation in NATO's history followed the September 11, 2001 attacks.\footnote{NATO Cooperative Cyber Defence Centre of Excellence. (n.d.). Cyber attacks and Article 5 – a note on a blurry but consistent position of NATO. \url{https://ccdcoe.org/incyder-articles/cyber-attacks-and-article-5-a-note-on-a-blurry-but-c} onsistent-position-of-nato/}

The stakes are concrete rather than theoretical:

\begin{itemize}
  \item More than 95 percent of internet traffic is carried by undersea cables
  \item Roughly 1.3 million kilometres of cabling underpins an estimated 10 trillion dollars of financial transactions every day
  \item A coordinated strike on the Estlink interconnectors or the regional data backbone is a direct hit to the energy and communications security of frontline member states.
  \item Each unanswered incident also lowers the cost of the next
\end{itemize}

The question in front of this committee is therefore not whether the threat is real, but how an alliance built to defend territory responds to a sustained campaign aimed at the seabed beneath it, using tools that international law was never written to provide.
